\chapter{Combinatorics}

\section{Three Principles of Combinatorics}

These principles are \emph{Addition}, \emph{Multiplication} and \emph{Inclusion-Exclusion}.

\section{The Basic Counting Principle}

For \(n\) outcomes of event \(A\) and \(m\) of \(B\), there are \(n \cdot m\) outcomes for event \(A\) and \(B\).

In combinatorics we can model a lot of problems as having \(n\) elements from which we can take one at a time 
a certain number of times, and then we take into consideration if we allow repetition or the order matters or not.

\section{Permutation Without Repetition}

The \emph{permutation} counts the number of ways of ordering \(n\) distinct elements. In other words, the number of tuples of \(n\) distinct elements.
Of course, the order matters.

\[
    n!
\]

\subsection{Permutation With Repetition}

The \emph{permutation with repetition} counts the number of ways of ordering \(n\) non-distinct elements. Here \(m_1, \dots, m_n\)
is the number of times a specific item is repeated in the original set. By dividing by the respective amount of the 
repeated non-distinct element, we eliminate the permutations where the identical elements swap their place.

\[
    \frac{n!}{m_1! \cdots m_n!}
\]

This formula comes from the \emph{multinomial coefficient} given by 

\[
    \binom{n}{n_1 n_2 \dots n_m} = \frac{n!}{n_1! \cdots n_m!},
\]

which comes from 

\[
    \binom{n}{n_1} \binom{n - n_1}{n_2}\binom{n - n_1 - n_2}{n_3} \dots \binom{n - n_1 - \cdots n_{m - 1}}{n_m}.
\]

This is also a kind of \emph{partition problem}, because it can be formulated as if you have 
\(n\) different elements, and you want to divide them into \(m\) bins with capacities \(m_1, m_2, \cdots, m_m\), how 
many possible ways are there for these partitions. 

In this case, the order matters, but only for the different elements. 

\textbf{Example:}

\section{Combination with Repetition}

We can think of the \emph{combination} as having \(n\) types of objects, and we want 
to take with repetition until we fill our \(k\) slots.

\[
    \binom{n + k - 1}{n - 1} 
\]

For understanding, we have \(n\) types of objects, and because we want a combination of those objects, the order does not matter.
Thus, we can put them together as contiguous groups. We have \(k\) slots and \(n\) types, thus we can separate our groups 
with \(n - 1\) separators \(|\). The question becomes how many ways are there to put \(n - 1\) separators into \(n + k - 1\) slots. Note 
that we have \(n + k - 1\) positions if we count the actual slots and the separator slots in between.

\[
    *|* * *|* *, \quad n = 6, \quad k = 2, \quad \text{positions in between } = 7
\]

\section{Variation}

The \emph{variation} counts the ways of putting \(n\) objects in \(k\) slots with repetition. Here, the order does not matter.

\[
    n^k
\]

\section{Variation Without Repetition}

This kind of \emph{variation}  counts the ways of putting \(n\) objects in \(k\) slots without repetition. Number of tuples with \(k\) 
elements which can be selected from \(n\) elements, without repetition. The order matters.

\[
    \frac{n!}{(n - k)!}
\]

\section{Combination Without Repetition}

The \emph{combination} can be understood as choosing \(k\) objects of \(n\) elements. Or putting \(k\) elements from \(n\) into \(k\) slots where the order does not matter.

\[
    \binom{n}{k} = \frac{n!}{k!(n-k)!}
\]

Note that like in the variation we use \(\frac{n!}{(n - k)!}\) to limit our slots to \(k\). Then for the \(k\) objects 
we choose we can arrange them in \(k!\) ways but because the order does not matter we divide by \(k!\) to get rid of the  
repeated combinations.

\section{Multinomial Coefficients}

When using the binomial coefficient repeatedly, one after the other, we can abbreviate 
these consecutive uses as 

\[
    \binom{n}{k_1, \dots, k_n}
\]

This is like from \(n\) choose \(k_1\), then from \(n - k_1\) choose \(k_2\), and so on.

\subsection{Multinomial Theorem}

Given \(x_1, \dots, x_p \in \R\) and \(\N\), then 

\[
    (x_1 + \cdots, x_n)^n = \sum_{k_1 + \cdots + k_{p = n}}  \binom{n}{k_1, \dots, k_p} x_{1}^{k_1} \cdots x_{p}^{k_p}
\]

\section{Example Problems}

\subsection*{Defected Machines}

A new delivery with 10 machines from which 2 are not working properly has to be tested. For that purpose a 
man in charge selects 5 of the machines randomly, and then he proceeds to make a decision based on the number of faulty machines found 
in the delivery. If among the 5 samples only 1 machine does not work it will be accepted otherwise it will be rejected. What is the 
probability of a delivery being accepted?

\textbf{Answer:}

Total cases: \(\binom{10}{5}\)

\(\prob(\text{accepted}) = \frac{\binom{8}{5}}{\binom{10}{5}} + \frac{\binom{2}{1}\binom{8}{4}}{\binom{10}{5}}\)

\(\binom{8}{5}\) represents no faulty machine in the choice; \(\binom{2}{1}\binom{8}{4}\) represents that only one machine does not work. 

\section{Disarray}

A \emph{disarray} is the number of permutations in which no object ends in the same initial spot. It is also known as the 
\emph{sub-factorial}.

\[
    !n = n! \sum_{k = 0}^{n} \left(\frac{{(-1)}^k}{k!}\right)
\]

\section{Bell Numbers}

The \emph{Bell numbers} count the number of ways of grouping \(n\) objects in an arbitrary number of slots.

\[
    B(N) = \sum_{k = 0}^{N-1}\binom{N}{K}B(K),
\]

with \(B(1) = 1\) and \(B(2) = 2\).

\section{Ramanujan Numbers}

The \emph{Ramanujan numbers} are the same as the Bell numbers, but all objects are equal. As an example, the number of ways
to decompose a number.

\[
    R(N) \approx \frac{1}{4N\sqrt{3}} e^{2\pi \sqrt{\frac{N}{6}}}
\]

\section{Stirling Numbers I}

The ways of permuting \(N\) items with \(K\) exchanges are given by the \emph{Stirling Numbers I}.

\[
    \genfrac{[}{]}{0pt}{}{N + 1}{K} = N \genfrac{[}{]}{0pt}{}{N}{K} + \genfrac{[}{]}{0pt}{}{N}{K - 1}
\]

\section{Stirling Numbers II}

The ways of grouping \(n\) items in \(k\) slots are given by the \emph{Stirling Numbers II}.

\[
    \genfrac{\{}{\}}{0pt}{}{N}{K} = \frac{1}{K!} \sum_{j=0}^{K} (-1)^{K - j} \binom{K}{j} j^N
\]

\section{Lah Numbers}

The \emph{Lah numbers} represent the ways of building \(K\) lists with a set of \(N\) elements.

\[
    \genfrac{\lfloor}{\rfloor}{0pt}{}{N}{K} = \binom{N - 1}{K - 1} \frac{N!}{K!} = \sum_{j=0}^{K} 
    \genfrac{[}{]}{0pt}{}{N}{K} \genfrac{\{}{\}}{0pt}{}{N}{K}
\]

\section{Euler's Numbers}

The \emph{Euler's numbers} are the number of permutations of a set of \(N\) elements of different size in which
\(K\) elements are bigger than their previous element.

\[
    \genfrac{\langle}{\rangle}{0pt}{}{N}{K} = \sum_{j = 0}^{K} (-1)^j \binom{N + 1}{j} (K + 1 - j)^N.
\]

\section{Catalan Numbers}

The \emph{Catalan numbers} have different applications like the number of ways of the triangulations of
a polygon or the number of trees with \(N\) leaves.

\[
    \mathfrak{C}(N) = \binom{2N}{N} - \binom{2N}{N + 1}
\]

or

\[
    \mathfrak{C}(N) = \sum_{k=1}^{n}(k - 1) \mathfrak{C}(N - k).
\]

